\documentclass[10pt]{article}
 
% Latin Modern: the modern successor to Computer Modern
\usepackage[T1]{fontenc}
\usepackage{lmodern}
\usepackage[utf8]{inputenc}
 
\usepackage[margin=0.5in]{geometry}
\usepackage{parskip}      % space between paragraphs, no indent
\usepackage{enumitem}
\usepackage{setspace}
\usepackage{multicol}
\singlespacing            % single-spaced body
 
\pagestyle{plain}
 

\begin{document}
 
% REMOVE THE PAGE NUMBER
\thispagestyle{empty}

% --- Header (top left) ---
\begin{multicols}{2}

    \noindent
    Daniel Kadyrov\\
    \today\\
    200-Hour Teacher Training 2026-2027\\
    Sonic Yoga NYC
 
% \vspace{1em}
    \columnbreak
    % align right
    \begin{flushright}
      Date: November 26th, 2026\\
      Day/Time: Thursday, 12:30 PM\\
      Class: Vinyasa 2 (60 min)\\
      Teacher: Mitch Hipp\\
    \end{flushright}
\end{multicols}


\vspace{1em}
\begin{center}
  \textbf{Yoga Class Observations}
\end{center}

\begin{enumerate}
    \item \textbf{Was there a theme or message for class? If so, how was it carried through?}\\
    This class was on the Thanksgiving holiday so the message of the class was giving thanks for yoga, Sonic, as well as us as practitioners. , 

    \item \textbf{Was there Pranayama, Meditation, Mantra or Mudra used?}\\
    Mitch cued both Ujjayi and Sama Vritti Pranayama. Throughout the class, Mitch also cued lions breath, Simha Pranayama. Mitch cued Kali mudra during crescent lunge, Anjaneyasana, and cued Anjali mudra during the twisted sequences, which included revolved lunge, Parivrtta Anjaneyasana, revolved side angle pose, Parivrtta Parsvakonasana, and twisted chair pose, Parivrtta Utkatasana. The mantra Ohm was initiated at the beginning and end of the class. Meditation was provided in the form of Shivasana.

    \item \textbf{Style of sequence: Cycle (waves of asana that repeat and build) or Category (Working on one category of postures at a time) or other?}\\
    The class follows a compartmentalized categorical format featuring cycling within some categories. Mitch splits sun salutations, standing poses, leg balances, twists, and floor poses with an Ashtanga style chataranga flow in between. Each categorical section features a cycle to krama the difficulty up and prepare for the peak pose that climaxes each partition. 
    
    \item \textbf{Was there a peak pose(s)?}\\
    The peak poses included the crow pose, Bakasana, the side crow pose, Parsva Bakasana, and a revolved side angle pose, Parivrtta Parsvakonasana.

    \item \textbf{Name the main actions of the peak pose(s):}\\
    Bakasana: Place hands on the mat, bend elbows, lift hips, and balance knees on upper arms. Engage core and gaze forward.\\
    Parsva Bakasana: Place hands on the mat, bend elbows, lift hips, and balance knees on upper arms. Twist torso to one side, engaging core and gazing forward.\\
    Parivrtta Parsvakonasana: From a high crescent lunge, twist torso towards the bent knee, bringing the opposite elbow to the outside of the bent knee. Optional bind with the opposite hand reaching behind the back.

    \item \textbf{What postures in the sequence prepared the body for the peak pose(s)?}\\
    Mitch prepped for Bakasana with a series of chataranga flows as well as placing the practioners in a low squat, Malasana. The side crow pose, Parsva Bakasana, was prepped with the earlier crow pose, Bakasana, with chair pose, Utkatasana, and with twists including revolved side angle pose, Parivrtta Parsvakonasana, and twisted chair pose, Parivrtta Utkatasana. The revolved side angle pose was prepped with the earlier revolved low crescent lunge, Parivrtta Anjaneyasana.

    \item \textbf{Give an example of pose/counterposes in the main sequence and in the finishing sequence.}\\
    In the standing sequence, Mitch cued triangle pose, Trikonasana, as a counter to the side angle pose, Parsvakonasana. After the leg balance, the tree pose, Vrksasana, Mitch cued a variation of a Shiva squat, Utkata Konasana, which stretched the opposite side of the body. During the seated sequence, Mitch alternated between forward folds, such as Bada Konasana and Paschimottanasana, with reverse tabletop, Ardha Purvottanasana, and reverse plank, Purvottanasana. 

    \item \textbf{What was the seated/finishing sequence?  Was it related to the standing sequence in any way or a counter energy sequence?}\\
    The final balancing sequence ended with a seated sequence that included a series of forward folds, such as Bada Konasana and Paschimottanasana, and reverse tabletop, Ardha Purvottanasana, and reverse plank, Purvottanasana. After a seated wide legged forward fold, Prasarita Padottanasana, Mitch cued a supine twist, Supta Matsyendrasana, and blissful baby, Ananda Balasana. Finally, Mitch cued the inversion, shoulder stand, Salamba Sarvangasana. As a counterpose and to respond to the high winds and inclimate weather outside, Mitch cued a flowing bridge sequence, Setu Bandha Sarvangasana, before closing the asana practice with Shivasana.

    \item \textbf{What did you learn from observing this class?}\\
    This was a rare occurrence where Mitch had more students tune in virtually than in person. As a result, Mitch had to adjust his cues and attention to accomodate the virtual students and keep them engaged. 

\end{enumerate}
\end{document}